% !TEX root = main_memo.tex
\chapter{Pointed Gluing and the General Structure}\label{sectionPointedGluing}

After \cref{SecondstepforBQOthm}, we now need to argue that \cref{Introthm:LevelsAreFinitelyGenerated} implies \cref{Introthm:BQO}.
The key to do so is to develop the technology to increase the $\CB$-rank\footnote{As we have seen, gluing functions does not.}, which is the objective of this section.
This is the motivation for introducing the Pointed Gluing operation from \cite{carroy2013quasi}, and study its binding conditions in the same way we did for the Gluing operation.
This will be central in \cref{PreciseStructureFinal,sectionDoubleSucc},
and gives the generalization to $\sC$ of \cite[Theorem 5.2]{carroy2013quasi} that we obtain at the end of the present section.

Note that the pointed gluing operation is implicitly used in the theory of Wadge degrees since Wadge's thesis~\cite{phdwadge}, sometimes with slight variations (see for instance \cite[Section 5]{selivanovnew}).
It has been used explicitly in~\cite[Subsection 2.6]{CMRSWadge}


Recall that for a sequence $s\in\N^{<\N}$ we denote by $s^n$ for $n\in\N$ the concatenation of $s$ with itself $n$-many times.
We make the convention that $s^0$ is the empty sequence and $s^\infty$ \index[IdxSymb]{s@{$s^n$, $s^\infty$}}is the infinite concatenation of $s$ with itself.


\begin{definition}
Given a sequence $(F_i)_{i\in \N}$ of subsets of $\cN$, 
%we call \emph{support} of $(F_i)_{i\in \N}$ the set of natural numbers $i$ for which $F_i\neq\emptyset$.
we call \emph{pointed gluing} of $(F_i)_ {i\in \N}$ the set $\pgl_{i\in\N} F_i =\{\iw{0}\}\cup\bigcup_{i\in \N}(0)^i\conc(1)\conc F_i$.

\index[IdxNotion]{operation!pointed gluing}\index[IdxSymb]{ptgl@{$\pgl_{i\in\N} F_i$}}\index[IdxSymb]{ptgl@{$\pgl_{i\in\N} f_i$}}
\medskip

Given a sequence of functions $(f_i:A_i\rao B_i)_{i\in \N}$ in $\functionsonbaire$,
we call \emph{pointed gluing} of $(f_i)_{i\in \N}$, the map
 \begin{align*}
\pgl_{i\in \N}f_i: \pgl_{i\in \N}A_i &\longrightarrow\pgl_{i\in \N}B_i\\
 x &\longmapsto\begin{cases}
(0)^i\conc(1)\conc f_i(x') & \mbox{if }x=(0)^i\conc(1)\conc x'\\
\iw{0} & \mbox{otherwise.}
\end{cases}
\end{align*}
\end{definition}

%\ynote{Do we really need the notion of support? For now the only reference in the paper, apart from here, is a footnote in \cref{ssec:centered_def}.}
%We call \emph{support} of a sequence $(f_i)_i$ of functions the support of $(\dom(f_i))_i$. If $(f_i)_i$ has empty support, then $\pgl(f_i)_i$ is the constant function $\iw{0}\mapsto \iw{0}$.

\begin{fact}\label{BasicsOnPointedGluing}
The pointed gluing operation preserves countability, closure, compactness, Polishness of sets; continuity, injectivity, surjectivity, scatteredness of functions.

It commutes with the identity natural transformation, i.e. $\id_{\pgl_{i}X_i}=\pgl_{i}\id_{X_i}$.

The point $\iw{0}$ is a continuity point of $\pgl_{i\in\N} f_i$ for any sequence of functions $(f_i)_{i\in\N}$ in $\functionsonbaire$.
\end{fact}

The following proposition generalizes \cite[Proposition 3.1]{carroy2013quasi}.

A sequence $(\alpha_n)_{n\in\N}$ of ordinals is called \emph{regular} when for all $m\in\N$ there exists a natural $n>m$ such that $\alpha_m\leq\alpha_n$, or
equivalently, when $\sup_n(\alpha_n+1)=\limsup_n(\alpha_n+1)$.

Recall that if $f$ is simple with $\CB(f)=\alpha+1$ the distinguished point\index[IdxNotion]{distinguished point (of a simple function)} of $f$ is the unique $y\in\im(f)$ satisfying $\set{y}=f(\CB_{\alpha}(f))$.

%\rnote{peut être avec le changement de $\sC$ on pourrait faire les rayons seulement dans $\cN$?}
%Given a $0$-dimensional separable metrizable space $B$ and a decreasing sequence $(U_n)_{n\in\N}$ of basic clopen neighborhood at $y\in B$, we define the \emph{$n$-th ray of $B$ at $y$} as the clopen set $\ray{B}{y,n}:=U_{n}\setminus U_{n+1}$. In such a case, if $f:A\to B$ is a function, the \emph{$n$-th ray of $f$ at $y$} is the function $\ray{f}{y,n}:=f\corestr{\ray{B}{y,n}}$.

%In general, these rays depend on the choice of a sequence of basic open neighborhoods for $y$,
%but if $B\subseteq\cN$ and $y\in B$ we always use the following a basic sequence of neighborhoods of $y$ given by $U_n=\nbhd{y}{n}\cap B$.

\index[IdxNotion]{Cantor--Bendixson!rank!of a pointed gluing}
\begin{proposition}\label{CBrankofPgluingofregularsequence1}
Let $f=\pgl_{n\in\N} f_n$ for some sequence $(f_n)_{n\in\N}$ of scattered functions in $\functionsonbaire$. 
If $(\CB(f_n))_{n\in\N}$ is regular and $\alpha=\sup_{n\in\N}(\CB(f_n))$, then $\CB_\alpha(f)=\set{\iw{0}}$.
In particular, $f$ is simple of distinguished point $\iw{0}$, and $\CB(f)=\alpha+1$.
\end{proposition}
\begin{proof}
Since $f_n\equiv f\restr{N_{(0)^n\conc(1)}}$, we have $\CB(f_n)=\CB(f\restr{N_{(0)^n\conc(1)}})$ by \cref{CBbasicsfromJSL-tp}.
If $\beta<\alpha$ then by regularity $\CB_\beta(f\restr{N_{(0)^n\conc(1)}})=\CB_\beta(f)\cap N_{(0)^n\conc(1)}$ is non-empty for infinitely many $n$,
which implies that $\iw{0}\in\CB_{\beta+1}(f)$. Therefore $\iw{0}\in\CB_{\alpha}(f)$.
Since $\CB_{\alpha}(f\restr{N_{(0)^n\conc(1)}})$ is empty for all $n$, it follows that $\CB_{\alpha}(f)=\set{\iw{0}}$.
\end{proof}

%\begin{aynote}{Remark: A simple function is not always the pointed gluing of its rays. Spoiler on the (infinitary) wedge operation.}
%\begin{example}\label{infiniteWedgeOfOnes}
%Consider the function defined by 
%\begin{align*}
%f: \gl_{n\in \N}\pgl_{m\in\N} \set{\iw{0}} &\to \pgl_{m\in \N} \gl_{n\in \N} \set{\iw{0}}\\
%(n)\conc (0)^m\conc (1)\conc \iw{0} &\mapsto (0)^m \conc  (1)\conc (n)\conc  \iw{0}\\
%(n)\conc \iw{0} &\mapsto \iw{0}
%\end{align*}
%Then $\CB(f)=2$ and $f$ is constant on $\CB_{1}(f)=\set{(n)\conc \iw{0}}[n\in\N]$, hence it is simple with distinguished point $y=\iw{0}$.
%For every $n\in\N$ we have $f\restr{N_{(n)}}\equiv \pgl_{m\in\N} \id_{\set{\iw{0}}}$, while $\CB(\ray{f}{y,m})\equiv \id_{\N}$ for all $m\in \N$. 
%It can be proved (see \cref{VerticalTheorem}) that $f\equiv  (\pgl_{m\in\N} \id_{\set{\iw{0}}}) \glbin \id_{\N}$ which lies strictly between $\pgl_{m\in\N} \id_{\set{\iw{0}}}$ and $\pgl_{m\in\N} \id_{\N}$, the pointed gluing of the rays of $f$.
%\end{example}
%\end{aynote}
%
%\begin{arnote}{au sujet de cet exemple}
%Une fois de plus, tu me vois un peu ennuyé: je suis convaincu du bien-fondé de cet exemple (à condition de dire explicitement ce que tu as mis dans la marge, tel quel il est difficile de voir où l'on veut en venir) mais je trouve que le mettre ici casse trop la lecture. C'est assez technique à lire, pour un point que beaucoup de lecteur pourront juger mineur, ou bien parfaitement évident... C'était le cas à Paris par exemple, intuitivement ils avaient tous bien en tête qu'il n'y a que très peu de chances qu'une fonction simple soit le recollement pointé de ses rayons.
%
%Par ailleurs, il y a déjà un exemple de ce gabarit dans le fait 7.8! En fait on a même montré dans la section 4 que ce sera le cas pour toute fonction simple non centrée...
%Je pense qu'on peut faire une remarque courte dans cette direction, et faire référence au fait 7.8, dans lequel on peut éventuellement rajouter l'exemple que tu as mis.
%Qu'en penses-tu? Yann: Je suis d'accord pour faire référence à une exemple dans la section 7 si tu penses que ça casse trop la lecture. 
%\end{arnote}

\begin{fact}\label{GluinglowerthanPgluing}
Given a sequence $(f_i)_{i\in \N}$ in $\functionsonbaire$,
we have $\gl_{i\in \N} f_i\leq \pgl_{i\in \N} f_i$.
\end{fact}

\begin{proof}
Use \cref{Gluingaslowerbound} with $f=\pgl_{i\in \N} f_i$ and $B_i=N_{(0)^i\conc(1)}$ for ${i\in \N}$.
\end{proof}

We write $x_n\to x$\index[IdxSymb]{_@{$x_n\to x$}} to say that a sequence $(x_n)_{n\in\N}$ of points in a topological space converges to a point $x$.
We recall the following useful sufficient criterion for continuity (\cite[Claim 3.2]{carroy2013quasi}).

\begin{lemma}\label{prop:sufficientcondforcont}
Let $A$ and $B$ be metrizable spaces and $f:A\rao B$.
Suppose that $U$ is an open subset of $A$ satisfying:
\begin{enumerate}
\item $f$ is continuous both on $U$ and on $A\setminus U$, and
\item for all $(x_n)_{n\in\N}$ in $U$, if $x_n\rao x\in A\setminus U$ then $f(x_n)\rao f(x)$.
\end{enumerate}
Then $f$ is continuous.
\end{lemma}

\begin{proof}
Since $A$ is metrizable, we prove sequential continuity at every point $x\in A$. Take a sequence $x_n\rao x$, we need to prove that $f(x_n)\rao f(x)$.
Observe first that if $x\in U$ then a tail of $(x_n)_n$ is actually included in the open set $U$, so using continuity of $f\restr{U}$ we have $f(x_n)\rao f(x)$.

So suppose that $x\notin U$, and partition $\N$ in two pieces: $I=\set{n\in\N}[x_n\in U]$ and $J=\N\setminus I$.
If $I$ is finite, then continuity of $f\restr{A\setminus U}$ yields $f(x_n)\rao f(x)$.
Otherwise, let $V\ni f(x)$ be an open set, the second hypothesis gives $N\in \N$ such that for all $n\in I$, if $n>N$ then $f(x_n)\in V$.
If $J$ is finite, fix $m=\max\set{N, \max J}$ to get $f(x_n)\in V$ for all natural numbers $n>m$.
If $J$ is infinite, use continuity of $f\restr{A\setminus U}$ to get $N'\in \N$ such that $f(x_n)\in V$ for all $n\in J$ with $n>N'$.
Therefore $f(x_n)\in V$ for all natural numbers $n>\max\set{N,N'}$, which proves that $f(x_n)\rao f(x)$ in this case too.
\end{proof}

\section{Pointed gluing as an upper bound}

We say that a sequence of functions $(f_i)_{i\in\N}$ \emph{is reducible by (finite) pieces to}\index[IdxNotion]{reducible by (finite) pieces} a sequence $(g_i)_{i\in\N}$ of functions
if there is a family $(I_n)_{n\in\N}$ of pairwise disjoint finite sets of $\N$ such that for all $n\in\N$ we have $f_n\leq\gl_{i\in I_n}g_i$.

For our analysis of the pointed gluing, we introduce the following way of partitioning a function at a given point of its codomain.

\begin{definition}
For $B\subseteq \cN$, $y\in \cN$ and $n\in \N$, the \emph{$n$-th ray of $B$ at $y$}\index[IdxNotion]{ray!of a set at a point} is the clopen subset of $B$ defined by 
\[
\ray{B}{y,n}:=\{x\in B\mid y\restr{n}\segm x \text{ and } y\restr{n+1} \nsegm x\}.
\]
In such a case, if $f:A\to B$ is a function, the \emph{$n$-th ray of $f$ at $y$}\index[IdxNotion]{ray!of a function at a point} is the function $\ray{f}{y,n}:=f\corestr{\ray{B}{y,n}}$.\index[IdxSymb]{B@{$\ray{B}{y,n}$}}\index[IdxSymb]{f@{$\ray{f}{y,n}$}}
\end{definition}




\index[IdxNotion]{upper bound condition!pointed gluing}
\index[IdxNotion]{pointed gluing operation!upper bound condition}
\begin{proposition}[Pointed gluing as upper bound]\label{Pgluingasupperbound}
Let $f\in \functionsonbaire$ be continuous %$A$ be a topological space, $B$ a 0-dimensional separable metrizable space
%$f:A\to B$ be a continuous function,
and $(g_i)_{i\in\N}$ be a sequence in $\functionsonbaire$.
If $y\in B$ and $(\ray{f}{y,j})_{j\in \N}$ is reducible by pieces to $(g_i)_{i\in \N}$,
then $f\leq\pgl_{i\in \N}g_i$.
\end{proposition}

\begin{proof}
Let $f:A\to B$ in $\functionsonbaire$ be continuous and note that if $y\notin\im f$ then we can actually conclude by \cref{Gluingasupperbound}, so suppose that $y\in\im f$.
Fix a family $(I_n)_{n\in\N}$ witnessing the reduction by pieces, and for all $n\in\N$ a reduction $(\sigma_n, \tau_n)$ from $\ray{f}{y,n}$ to $\gl_{i\in I_n}g_i$.
Observe that for all $n\in\N$ we have $\sigma_n(x)=(i)\conc x'$ for some $i\in I_n$ and $x'\in \dom g_i$.
Since $f$ is continuous and the sets $\ray{B}{y,n}$ are clopen, the sets $A_n:=\dom(\ray{f}{y,n})=\dom(\sigma_n)$ form a clopen partition of $A\setminus f^{-1}(\set{y})$,
so $\tilde{\sigma}=\bigsqcup_n\sigma_n$ is a well-defined function on $A\setminus f^{-1}(\set{y})$.
Define $\sigma$ as follows: $\sigma(x)=\iw{0}$ if $f(x)=y$, and $\sigma(x)=(0)^i\conc(1)\conc x'$ when $\tilde{\sigma}(x)=(i)\conc x'$.

Let $I=\bigsqcup_{n\in\N} I_n$ and $\eta:I\to \N$ be such that $\eta(i)$ is the unique $n\in\N$ such that $i\in I_n$.
Define now $\tau$ as follows: $\tau(\iw{0})=y$, and $\tau(x)=\tau_{\eta(i)}((i)\conc x')$ when $x=(0)^i\conc(1)\conc x'$, $(i)\conc x'\in \im \tau_{\eta(i)}$ and $i\in I$.

As the sets $I_n$ are pairwise disjoint, $\eta$ is injective and $\tau$ is well-defined.
Note that, by construction, $\sigma$ is defined on all of $A$ and $f=\tau(\pgl_{i\in \N}g_i)\sigma$.
So it suffices to show that $\sigma$ and $\tau$ are continuous.

\medskip

We now use \cref{prop:sufficientcondforcont} with the open set $U=\dom f \setminus f^{-1}(\set{y})=\bigcup_{n\in\N}A_n$ to obtain continuity of $\sigma$:
as $\tilde{\sigma}$ is continuous by \cref{lem:ContUnion} so is $\sigma$ on $U$, hence it is enough to check the second condition.
Towards that goal, suppose that a sequence $(x_n)_n$ in $\dom f$ with $f(x_n)\neq y$ converges to $x$ with $f(x)=y$. To show that $\sigma(x_n)\rao\iw{0}$, we prove that given any $m\in\N$, we have $\sigma(x_n)\in N_{(0)^m}$ for sufficiently large $n$.

For all $n\in\N$, let $j_n\in\N$ be such that $f(x_n)\in \ray{B}{y,j_n}$. By definition of $\sigma$ we have $\sigma(x_n)\in N_{(0)^i}$ for some $i\in I_{j_n}$.
Now, as the sets $(I_{j_n})_{n\in\N}$ are nonempty and pairwise disjoint, we have $\min I_{j_n}\to\infty$ as $n\to \infty$.
So there exists $N\in \N$ such that for all $n\geq N$ we have $\min I_{j_n}\geq m$ and therefore $\sigma(x_n)\in N_{(0)^m}$.

We proceed similarly for $\tau$ using this time $U=\dom \tau \setminus\set{\iw{0}}$. Note that $\tau$ is continuous on $U$ by \cref{lem:ContUnion}. 
Now take a sequence $(x_n)_{n\in\N}$ in $U$ converging to $\iw{0}$, we prove that $\tau(x_n)\rao y$. %Since $I\to \N$, $i\mapsto n_i$ 
Given $m\in\N$, set $M=\max(\bigcup_{j\leq m}I_j)$. Since $x_n\rao\iw{0}$ there is $N\in\N$ such that $x_n\in N_{(0)^{M+1}}$ for all $n>N$.
Therefore by definition of $\tau$, for all $n>N$ we have $\tau(x_n)\in\bigcup_{j>m}\im(\tau_j)$, which means that $\tau(x_n)\in \nbhd{y}{m}$.
We thus get $\tau(x_n)\rao y$, so $\tau$ is continuous by \cref{prop:sufficientcondforcont}.
\end{proof}

\cref{Pgluingasupperbound} gives the following rough -- yet very useful -- way to upper bound a continuous function using the pointed gluing operation.

\begin{corollary}\label{Pgluingofraysasupperbound}
Let $f:A\to B$ be a continuous function in $\functionsonbaire$.
Then $f\leq \pgl_{i\in \N}\ray{f}{y,i}$ for all $y\in B$.
\end{corollary}

%\ynote*{support again... do we really need this?}{
%To simplify notation, when $I\subseteq \N$ and $(g_i)_{i\in I}$ is a sequence of functions indexed by $I$, we adopt the convention that $\pgl_{i\in I}g_i:=\pgl_{i\in\N}f_i$,
%where $f_i=g_i$ for all $i\in I$ and $f_i=\emptyset$ otherwise.}

The following corollary states another relationship of the pointed gluing with the finite gluing:

\begin{corollary}\label{SplittingaPgluingonatail}
Let $(f_i)_{i\in \N}$ be continuous in $\functionsonbaire$.
For all $n\in\N$, we have $\pgl_{i\in\N}f_i\equiv(\gl_{i<n}f_i)\glbin \pgl_{i\geq n}f_{i}$.
\end{corollary}
\begin{proof}
Using \cref{Pgluingasupperbound} with $y=\iw{0}$, we get $(\gl_{i<n}f_i)\glbin \pgl_{i\geq n}f_{i}  \leq \pgl_{i\in \N} f_i$.
Moreover, we see that  $\pgl_i f_i\leq (\gl_{i<n}f_i)\glbin \pgl_{i\geq n}f_{i}$ using \cref{Gluingaslowerbound} with the clopen partition $\set{N_{(0)^i\conc(1)}}[i<n]\cup\set{N_{(0)^n}}$.
\end{proof}

Recall that we saw in~\cref{CBrankofPgluingofregularsequence1} that the pointed gluing of a sequence with regular $\CB$-ranks is simple. Given a simple function $f$ with distinguished point $y$, we have $f\leq \pgl_{n\in\N} \ray{f}{y,n}$ by \cref{Pgluingofraysasupperbound}. While it is not true in general (see \cref{SomeCounterExamples}) that $f\equiv \pgl_{n\in\N} \ray{f}{y,n}$, we have the following:

\begin{proposition}\label{CBrankofPgluingofregularsequence2simple}
If $f\in\functionsonbaire$ is scattered of $\CB$-rank $\alpha+1$ is simple of distinguished point $y$, then the sequence $(\CB(\ray{f}{y,n}))_{n\in\N}$ is regular and its supremum is $\alpha$. 
\end{proposition}

\begin{proof}
Take a simple $f:A\to B$ is in $\sC_{\alpha+1}$ for some $\alpha<\omega_1$ with distinguished point $y$, and set $\alpha_i=\CB(\ray{f}{y,i})$ for all $i\in\N$.
By simplicity, $\CB_\alpha(f)\subseteq f^{-1}(\{y\})$ and so for all $i\in \N$ we have $\CB_\alpha(\ray{f}{y,i})=\CB_\alpha(f)\cap f^{-1}(\ray{B}{y,i})=\emptyset$,
which implies $\alpha_i\leq \alpha$. Hence $\sup_i\alpha_i \leq \alpha$.
To get the other inequality as well as regularity,
we need to show that for all $\beta<\alpha$ and for all $m\in\N$ there is an $n>m$ such that $\beta<\alpha_n$. 

Suppose towards a contradiction that for some $\beta<\alpha$ for all $n>m$ we have $\alpha_n\leq\beta$. Setting $g:=f\restr{\dom(f)\setminus(\bigcup_{i\leq m}\dom(\ray{f}{y,i}))}$
we observe that $\CB_\beta(g)\subseteq f^{-1}(\set{y})$ because $\CB_\beta(\ray{g}{y,n})$ is always empty: for $n\leq m$ it is by definition of $g$, and for
$n>m$ this follows from our assumption that $\ray{g}{y,n}=\ray{f}{y,n}$ has $\CB$-rank $\alpha_n\leq\beta$ using \cref{CBbasics0}~\cref{CBbasicsfromJSL2}.
Thus $g\restr{\CB_\beta(g)}$ is either empty or constant, in any case $\CB(g)\leq\beta+1\leq\alpha$.
Observe finally that $f=g\sqcup(\bigsqcup_{i\leq m}\ray{f}{y,i})$ is a disjoint union of functions of $\CB$-rank $\leq\alpha$, so by \cref{CBrankofclopenunion}
we get $\CB(f)\leq\alpha<\alpha+1$, a contradiction.
\end{proof}


\section{Maximum functions}

As an important consequence of the ``upper bound'' criterion stated in \cref{Pgluingasupperbound}, we identify functions which are maximum among functions in $\sC$ with $\CB$-rank at most $\alpha$.

Recall that $\omega f$\index[IdxSymb]{w@{$\omega f$}} denotes the gluing of the constant sequence with value $f$. Similarly, we write $\pgl f$ for the pointed gluing of the constant sequence with value $f$.

\begin{definition}[Minimum and maximum functions]\label{Def_MinMaxFunc}
We define by induction on $\alpha$ a function \index[IdxNotion]{minimum and maximum functions}$\Maximalfct{\alpha}$ in $\sC_\alpha$ and a function $\Minimalfct{\alpha+1}$ in $\sC_{\alpha+1}$.
Set $\Maximalfct{0}=\emptyset$ and $\Minimalfct{1}=\pgl\emptyset$.\index[IdxSymb]{l@{$\Maximalfct{\alpha}$}}\index[IdxSymb]{k@{$\Minimalfct{\alpha+1}$}}
Suppose that for all $\beta<\alpha$ the functions $\Maximalfct{\beta}$ and $\Minimalfct{\beta+1}$ are defined.

If $\alpha=\beta+1$ for some $\beta<\alpha$, then set $\Minimalfct{\alpha+1}=\pgl \Minimalfct{\alpha}$ and $\Maximalfct{\alpha}=\omega\pgl \Maximalfct{\beta}$.

Otherwise, that is if $\alpha$ is limit, then fix a sequence $(\alpha_n)_{n\in\N}$ cofinal in $\alpha$ and set $\Minimalfct{\alpha+1}=\pgl_n\Minimalfct{\alpha_n+1}$.
Fix as well an enumeration $(\beta_n)_{n\in\N}$ of $\alpha$ and set $\Maximalfct{\alpha}=\gl_n\Maximalfct{\beta_n}$.
\end{definition}

Note that $\Minimalfct{1}\equiv\id_{1}$ and $\Maximalfct{1}\equiv \id_{\N}$. For $\alpha$ limit, up to continuous equivalence, the definition of these functions does not depend on the choice of the sequences $(\alpha_n)_n$ and $(\beta_n)_n$.
For $\Minimalfct{\alpha+1}$, it follows from a classical well-known result of Mazurkiewicz and Sierpinskì \cite[Théorème 1]{SierpMazCompDen}
- we will get an alternative proof of that later with \cref{Compactdomains};
as for $\Maximalfct{\alpha}$, by \cref{Gluingasupperbound,Gluingaslowerbound},
different choices of enumerations $(\beta_n)_n$ produce continuously equivalent functions $\Maximalfct{\alpha}$.

Observe that since the Gluing and Pointed Gluing operations commute with the identity natural transformation and $\Minimalfct{1}=\id_{\set{\iw{0}}}$, the functions from \cref{Def_MinMaxFunc} are identity functions on their domains.
The spaces $\dom \Minimalfct{\alpha}$ are somewhat classical objects, they were defined exactly this way in \cite{carroy2013quasi}.
Their properties can prove useful in other -- yet related -- contexts, as illustrated by \cite[Proposition 3.12]{CMRSWadge}.
We will see here (in \cref{Minfunctions}) that each function $\Minimalfct{\alpha+1}$ is minimum in $\sC_{\geq \alpha+1}$.


The functions $\Maximalfct{\alpha}$ were defined in \cite[Definition 5.34]{phdcarroy}
and a version of the following result was proved for functions with Polish $0$-dimensional domain (\cite[Proposition 5.36 and Corollary 6.6]{phdcarroy}).


\begin{Proposition}[Maximum functions]\label{Maxfunctions} \index[IdxNotion]{maximum function proposition}
For all $\alpha<\omega_1$:
\begin{enumerate}
\item the function $\Maximalfct{\alpha}$ is a maximum of $\sC_{\leq\alpha}$,
\item the function $\pgl \Maximalfct{\alpha}$ is a maximum for simple functions in $\sC_{\leq \alpha+1}$,
\item for all $n\in\N$, $(n+1)\Minimalfct{\alpha+1}$ is a maximum among functions of $\CB$-type $(\alpha+1,n+1)$ with compact domains.
\end{enumerate}
\end{Proposition}

\begin{proof}
First notice that if $\alpha\leq \beta$, then we have $\Maximalfct{\alpha}\leq\Maximalfct{\beta}$ and $\Minimalfct{\alpha+1}\leq\Minimalfct{\beta+1}$.
For $\alpha=0$, we have $\Maximalfct{0}=\emptyset$ and $\Minimalfct{1}=\pgl\Maximalfct{0}=\id_{\set{\iw{0}}}\equiv\id_{\set{0}}$, so all items follows from \cref{LocallyConstantFunctions}.

\smallskip

We prove the first two items simultaneously by strong induction on $\alpha$: suppose they both hold for all $\beta<\alpha$. 
To see that $\Maximalfct{\alpha}$ is a maximum in $\sC_{\leq\alpha}$, let $f\in\sC$ with $\CB(f)\leq\alpha$.
By the Decomposition \cref{JSLdecompositionlemma}, $f$ is locally simple. If $\alpha$, is limit $f=\bigsqcup_i f_i$ with $f_i$ simple and $\CB(f_i)=\beta_i+1<\alpha$ and so by induction hypothesis the second item ensures that $f_i\leq \pgl \Maximalfct{\beta_i}\leq \Maximalfct{\beta_i+1}$. If $\alpha$ is successor, $f=\bigsqcup_i f_i$ with $f_i$ simple and $\CB(f_i)=\beta+1=\alpha$ and again the induction hypothesis implies that $f_i\leq \pgl \Maximalfct{\beta}$. In both, cases we get $\gl_{i}f_i\leq \Maximalfct{\alpha}$ and so $f\leq \Maximalfct{\alpha}$ by \cref{Gluingasupperbound}.

Now take $f$ simple in $\sC_{\leq \alpha+1}$ and call $y$ its distinguished point.
By \cref{Pgluingofraysasupperbound} we have $f\leq\pgl_{j\in\N}\ray{f}{y,j}$, but by simplicity of $f$ we also have $\CB(\ray{f}{y,j})\leq\alpha$ for all $j\in\N$. As $\Maximalfct{\alpha}$ is a maximum in $\sC_{\alpha}$, we get $\ray{f}{y,j}\leq\Maximalfct{\alpha}$ for all $j\in \N$ and so $f\leq \pgl\Maximalfct{\alpha}$ by \cref{Pgluingasupperbound}.

\smallskip

We finally prove the second point by a double induction. Suppose that for all $\beta<\alpha$, if $f\in \sC$ has compact domain and $\CB$-type $(\beta+1,n+1)$
for some natural $n$, then $f\leq(n+1)\Minimalfct{\beta+1}$.

Take now $f$ with compact domain and $\CB(f)=\alpha+1$. Suppose first that $f$ is simple, that is $n=0$.
Then for all $j\in \N$ the $\CB$-rank of the $\ray{f}{y,j}$ is at most $\alpha$ by \cref{CBrankofPgluingofregularsequence2simple}.
Moreover, each ray has compact domain, so it is either empty or $\tp(\ray{f}{y,j})=(\beta_j+1,n_j+1)$ for some $\beta_j<\alpha$ and $n_j\in\N$ by \cref{CBbasics0}.
In any case, by induction hypothesis the sequence $(\ray{f}{y,j})_j$ is reducible by pieces to the sequence
$(\Minimalfct{\alpha_n})_{n\in\N}$, where $(\alpha_n)_n$ is either constant if $\alpha$ is itself successor, or cofinal in $\alpha$ if it is limit.
In any case, by \cref{Pgluingasupperbound} we obtain $f\leq\Minimalfct{\alpha+1}$, as desired.

Suppose now that $tp(f)=(\alpha+1,n+1)$ for some $n>0$, then by \cref{FiniteDegreeAreFinGl} we have $f\equiv\gl_{i\leq n}f_i$ with $f_i$ simple for all $i\leq n$.
Now using the simple case, we indeed get $f\leq (n+1)\Minimalfct{\alpha+1}$.
\end{proof}



\section{Pointed gluing as a lower bound and minimum functions}

Given a sequence $(A_n)_{n\in\N}$ of subsets of a space $A$, and a point $a\in A$, we say that $(A_n)_n$ \emph{converges} to $a$ and we write $A_i\rao a$\index[IdxSymb]{__@{$A_i\rao a$}}
when for any open neighborhood $U$ of $a$, there is a natural number $m$ such that for all $n\geq m$ we have $A_n\subseteq U$. For instance, in any space $B$, the sequence of rays $(\ray{B}{y,j})_{j\in \N}$ at some point $y$ converges to $y$.

Characterizing when a pointed gluing reduces continuously to a given function appears to be considerably harder. The following rough criterion is sufficient for the sequel.


\begin{lemma}\label{Pgluingaslowerbound}
Let $f:A\to B$ be a function between metrizable spaces and $(g_n)_{n\in\N}$ be a sequence in $\functionsonbaire$.

Assume that there is a point $x\in A$ and a sequence $(A_n)_{n\in\N}$ of clopen sets of $A$ satisfying:
\begin{enumerate}
\item For all $n\in\N$, $f(x)\notin \overline{f(A_n)}$, 
\item The sets $f(A_n)$, for $n\in\N$, form a relative clopen partition, 
\item $A_n\rao x$,
\item For all $n\in\N$, we have $g_n\leq f\restr{A_n}$.
\end{enumerate}
Then $\pgl_{n\in \N}g_n\leq f$.
\end{lemma}

\begin{proof}
Choose for every $n\in\N$ a reduction $(\sigma_n, \tau_n)$ from $g_n$ to $f\restr{A_n}$.
We define $\sigma$ and $\tau$ as follows: set $\sigma(\iw{0})=x$, $\sigma((0)^n\conc(1)\conc x')=\sigma_{n}(x')$ if $x'\in \dom \sigma_n$, and $\tau(f(x))=\iw{0}$,
$\tau(y)=(0)^n\conc(1)\conc\tau_n(y)$ if $y\in \dom(\tau_n)$. Note that $\tau$ is well-defined since the sets $\dom(\tau_n)\subseteq f(A_n)$ are pairwise disjoint and do not contain $f(x)$.

The pair $(\sigma,\tau)$ is the desired reduction, it remains to show that it is continuous.
To see that $\sigma$ is continuous, we proceed using \cref{prop:sufficientcondforcont} with $U=\dom \sigma \setminus \{\iw{0}\}$. %\bigcup_{n\in\N}N_{(0)^n\conc(1)}$
%to show that $\sigma$ is continuous. 
As $\sigma$ is continuous on $U$ by \cref{lem:ContUnion}, we check the second condition of \cref{prop:sufficientcondforcont}.
Any sequence $(x_n)_n$ in $\dom \sigma \setminus \{\iw{0}\}$ converging to $\iw{0}$ yields by definition
a sequence $(k_n)_n$ in $\N$ such that $k_n\to \infty$ and $x_n\in N_{(0)^{k_n}\conc(1)}$ for all $n\in\N$.
Hence $\sigma(x_n)\in A_{k_n}$ for all $n$ by definition of $\sigma$ and since $A_n\rao x$, it follows that $\sigma(x_n)\rao x$, which proves that $\sigma$ is continuous.

To prove the continuity of $\tau$ we use \cref{prop:sufficientcondforcont} again, this time with $U=\dom \tau \setminus \set{f(x)}$. %$\bigcup_{n\in\N}f(A_n)$.
As the sets $f(A_n)$ form a relative clopen partition, so do the sets $\dom \tau_n$, so $\tau$ is continuous on $U$ by \cref{lem:ContUnion}. To prove the second condition, take a sequence $(y_n)_n$ in $\im (f\sigma) \setminus\set{f(x)}$ such that $y_n\rao f(x)$.
For all $n$, let $k_n$ be the unique natural number such that $y_n\in f(A_{k_n})$. We need to show that $\tau(y_n)\rao\iw{0}$, and therefore, by definition of $\tau$, it is enough to see that $k_n\rao\infty$.
But this is indeed the case, since otherwise $(y_n)_n$ would admit a subsequence included in a set $f(A_i)$ for some $i\in\N$
and therefore we would have $f(x)\in \overline{f(A_i)}$, a contradiction.
\end{proof}

The conditions of the previous statement are quite difficult to satisfy in practice. In fact, we exclusively make use of the following specific case which relies on conditions that are simpler to verify, yet apparently much stronger. % the following restrictive situation.
\index[IdxNotion]{lower bound condition!pointed gluing}
\index[IdxNotion]{pointed gluing operation!lower bound condition}
\begin{proposition}[Pointed gluing as lower bound]\label{Pgluingaslowerbound2}
Let $f:A\to B$ be a continuous function in $\functionsonbaire$ and $(g_i)_{i\in\N}$ be a sequence in $\functionsonbaire$.

Assume that %there is a point $x\in A$ such that $f$ is continuous at $x$ and
for all $i\in\N$ and all open neighborhood $U\ni x$
there is a continuous reduction $(\sigma,\tau)$ from $g_i$ to $f$ such that $\im (\sigma )\subseteq U$ and $f(x)\notin\overline{\im(f\sigma)}$.
Then $\pgl_{i\in \N}g_i\leq f$.

In fact, $\pgl_{i\in\N}g_i\leq f\restr{V}$ for all clopen neighborhood $V$ of $x$.
\end{proposition}

\begin{proof}
We define by induction a sequence $(A_n)_n$ of clopen subsets of $A$ such that the sets $f(A_n)$ are separated by open sets,
for all $n\in\N$ we have $f(x)\notin \overline{f(A_n)}$, $A_n\subseteq \nbhd{x}{n}$ and there exists a continuous reduction $(\sigma_n,\tau_n)$ from $g_n$ to $f\restr{A_n}$.
This sequence satisfies the conditions of \cref{Pgluingaslowerbound}, so the result follows.


Assume that the sets $(A_i)_{i<n}$ have been defined as above.
Let $k$ be such that $N_{f(x)\restr{k}}$ is disjoint from the closed set $\bigcup_{i<n}\overline{f(A_i)}$.
Use continuity of $f$ at $x$ to choose a clopen neighborhood $U$ of $x$ included in $N_{x\restr{n}}$ such that $f(U)\subseteq N_{f(x)\restr{k}}$.
The hypothesis guarantees the existence of a continuous reduction $(\sigma_n,\tau_n)$ from $g_n$ to $f\restr{U}$ with $f(x)\notin\overline{\im(f\sigma_n)}$.
This allows us to choose a clopen neighborhood $W\ni f(x)$ disjoint from $\overline{\im(f\sigma_n)}$ and we set $A_n=U\setminus f^{-1}(W)$.
The set $A_n\subseteq U\subseteq N_{x\restr{n}}$ is clopen as a difference of two clopen sets, $f(A_n)\cap W=\emptyset$ so $f(x)\notin \overline{f(A_n)}$,
$f(A_n)\subseteq N_{f(x)\restr{k}}\setminus W$ so $N_{f(x)\restr{k}}\setminus W$ is an open set witnessing separation of the sets $A_n$.
Finally, notice that $(\sigma_n,\tau_n)$ is actually a reduction from $g_n$ to $f\restr{A_n}$ since $\im \sigma_n \subseteq A_n$.
 \end{proof}

Seeing pointed gluing as an upper bound gave us maximum functions in $\sC_{\leq \alpha}$ (\cref{Maxfunctions}); similarly the lower bound criterion gives us minimum functions in $\sC_{\geq \alpha+1}$. %
%of the functions $\Minimalfct{\alpha+1}$.

\begin{proposition}[Minimum functions]\label{Minfunctions}\index[IdxNotion]{minimum function proposition}
For all $\alpha<\omega_1$ and $f\in \sC$, if $\alpha+1\leq \CB(f)$, then $\Minimalfct{\alpha+1}\leq f$.
\end{proposition}

\begin{proof}
We prove by strong induction for every $\alpha<\omega_1$, the following statement: $\Minimalfct{\alpha+1}\leq f$ for every simple function $f\in \sC_{\alpha+1}$. To see that this is enough, we remark that if this statement holds for all $\beta<\alpha$, then actually for all $f\in \sC_{<\alpha}$ if $\beta+1\leq\CB(f)$ then $\Minimalfct{\beta+1}\leq f$. 

To prove this remark, suppose that  $f\in \sC$ with $\CB(f)<\alpha$. Since $f$ is locally simple by the Decomposition \cref{JSLdecompositionlemma}, we can write $f=\bigsqcup_{i\in I} f_i$ with $f_i$ simple, and $\sup_{i\in I} \CB(f_i)=\CB(f)$ by \cref{CBrankofclopenunion}. So if $\beta+1\leq \CB(f)$, $\beta+1\leq \CB(f_i)=\beta_i+1$ for some $i\in I$ and we have $\Minimalfct{\beta+1}\leq \Minimalfct{\beta_i+1} \leq f_i\leq f$, where $ \Minimalfct{\beta_i+1} \leq f_i$ is granted by the statement, which proves the remark. 


For $\alpha=0$, note that $\Minimalfct{1}$ continuously reduces to any non\-empty function.
So suppose that $\alpha>0$ and that the statement holds for every $\beta<\alpha$.

Take a simple function $f\in\sC_{\alpha+1}$ and let $y\in\im f$ be the distinguished point of $f$.
Seeing that $\Minimalfct{\alpha+1}$ is a pointed gluing, we seek to apply \cref{Pgluingaslowerbound2} for some point $x\in\CB_\alpha(f)$. Let $U\ni x$ be any open set of $\dom f$. 
Notice first that as $x\in\CB_\alpha(f)$ and $x\in U$ we get that $f_U=f\restr{U}$ is simple and $\CB(f_U)=\alpha+1$ by \cref{CBbasics0}~\cref{CBbasicsfromJSL2}. 
Note that by \cref{CBrankofPgluingofregularsequence2simple} the sequence $(\CB(\ray{f_{U}}{y,n}))_n$ is regular of supremum $\alpha$. If $\alpha=\beta+1$ is successor, then there exists $N$ such that $\alpha=\CB(\ray{f_U}{y,N})$. By the induction hypothesis combined with our first remark, we get $\Minimalfct{\alpha}\leq \ray{f_{U}}{y,N}$.
Notice that if $(\sigma, \tau)$ witnesses this reduction then both $\im \sigma \subseteq U$ and $\overline{\im f\sigma}\subseteq \ray{B}{y,N}$ hold. This ensures that $\Minimalfct{\alpha+1}=\pgl \Minimalfct{\alpha} \leq f$ by \cref{Pgluingaslowerbound2}.


If $\alpha$ is limit, for all $\beta< \alpha$ there exists $N$ such that $\beta+1\leq \CB(\ray{f_{U}}{y,N})< \alpha$. So by the induction hypothesis combined with our first remark, we get $\Minimalfct{\beta+1}\leq \ray{f_{U}}{y,N}$.
So similarly as in the successor case, we get $\Minimalfct{\alpha+1}=\pgl_{k}\Minimalfct{\beta_k+1}\leq f$, for $(\beta_k)_k$ cofinal in $\alpha$.
\end{proof}



\section{General Structure}

We can now complete the analysis of all functions in $\sC$ with compact domain, thus generalizing \cite[Theorem 4.2]{carroy2013quasi}.

\index[IdxNotion]{functions with compact domain (classification)}
\begin{theorem}[Classification of functions with compact domains]\label{Compactdomains}
Assume that $f$ and $g$ are two functions in $\sC$ with a compact domain, then $f$ continuously reduces to $g$ if and only if $tp(f)\leq_{lex}tp(g)$.
More specifically $f\equiv(n+1)\Minimalfct{\alpha+1}$, where $tp(f)=(\alpha+1,n+1)$.

In particular, continuous reducibility is a pre-well-order\footnote{A \emph{pre-well-order} is a linear \wqo{} (for more on this see \cite[Section 34.A]{kechris}).}
of length $\omega_1$ on functions in $\sC$ with compact domain.
\end{theorem}

\begin{proof}
By definition for $\alpha<\beta$ we have $\Minimalfct{\alpha+1}\leq\Minimalfct{\beta+1}$ and since $\CB(\Minimalfct{\alpha+1})=\alpha+1<\beta+1=\CB(\Minimalfct{\beta+1})$
we get $\Minimalfct{\beta+1}\not\leq\Minimalfct{\alpha+1}$, so the functions $\Minimalfct{\alpha+1}$ for $\alpha<\omega_1$ form a well-order of length $\omega_1$.

Therefore by \cref{Gluingcohomomorphism} it is enough to prove that for $f\in\sC$ with compact domain then $f\equiv(n+1)\Minimalfct{\alpha+1}$, where $tp(f)=(\alpha+1,n+1)$.
Now by \cref{FiniteDegreeAreFinGl} we only have to show it for simple functions, which follows from \cref{Maxfunctions,Minfunctions}.
\end{proof}

In particular we obtain the promised alternative proof that the definition of $\Minimalfct{\alpha+1}$ does not depend on the choice of the cofinal sequence when $\alpha$ is limit.

The results on maximum and minimum functions can be leveraged to obtain important information about the structure of continuous reducibility on the whole of $\sC$.
As announced, this generalizes \cite[Theorem 5.2]{carroy2013quasi} and we (still) refer to it as the General Structure Theorem.\index[IdxNotion]{General Structure Theorem}

\input{general_structure_diagram_small.tex}

\begin{theorem}[General Structure]\label{JSLgeneralstructure}
For all functions $f$ and $g$ in $\sC$, and all ordinals $\lambda<\omega_1$ limit or null, we have:
\begin{enumerate}
\item If $\CB(f)\leq \CB(g)=\lambda$, then $f\leq g$.
\item For all $n\in \N$, $\CB(f)=\lambda+n$ and $\lambda+2n+1\leq \CB(g)$ imply $f\leq g$.
\end{enumerate}
In particular, $2\CB(f)<\CB(g)$ implies $f\leq g$.
\end{theorem}

\begin{proof}
Proceed by induction on $\lambda<\omega_1$. Suppose that the theorem is proven for all $\alpha<\lambda$ with $\lambda$ limit or null.

\smallskip

We first prove the second item. As $f\leq\Maximalfct{\lambda+n}$ by \cref{Maxfunctions} and $\Minimalfct{\lambda+2n+1}\leq g$ by \cref{Minfunctions}, it is enough to prove that
$\Maximalfct{\lambda+n}\leq \Minimalfct{\lambda+2n+1}$ for all $n\in\N$, and to do so we proceed by induction on $n\in\N$.
For $n=0$, if $\lambda=0$ then $\Maximalfct{0}=\emptyset\leq\Minimalfct{1}$, so we can suppose that $\lambda\neq0$.
Take $(\alpha_n)_n$ cofinal in $\lambda$ and $(\beta_n)_n$ an enumeration of $\lambda$ then, by induction hypothesis, for some injection $p:\N\rao\N$ we have
$\Maximalfct{\alpha_n}\leq\Minimalfct{\beta_{p(n)}+1}$ so by \cref{Gluingasupperbound,GluinglowerthanPgluing} we get
\(\Maximalfct{\lambda}\equiv\gl_n\Maximalfct{\alpha_n}\leq\gl_n\Minimalfct{\beta_n+1}\leq\pgl_n\Minimalfct{\beta_n+1}\equiv\Minimalfct{\lambda+1}.\)

If now $\alpha=\lambda+(n+1)$ then using the induction hypothesis, \cref{Pgluingasupperbound,GluinglowerthanPgluing} we see that
$\Maximalfct{\alpha}\equiv\omega\pgl\Maximalfct{\lambda+n}\leq\pgl\pgl\Minimalfct{\lambda+2n+1}\equiv\Minimalfct{\lambda+2n+3}= \Minimalfct{\lambda+2(n+1)+1}$.

\smallskip

To see the first item, take a function $g:A\to B$ of $\CB$-rank $\lambda$. By \cref{Maxfunctions} $f\leq\Maximalfct{\lambda}$ for any $f\in \sC_{\leq \alpha}$, so it is enough to show that if $\Maximalfct{\lambda}\leq g$.
If $\lambda=0$ then $g$ is the empty function, so suppose that $\lambda$ is limit. 

We are going to find a sequence $(s_n)_{n\in\N}\subseteq\N^{<\N}$ of finite sequences pairwise incomparable for the prefix relation such that the sequence $(\CB(g\corestr{N_{s_n}}))_n$
is either constant equal to $\lambda$ or strictly below $\lambda$ and cofinal in $\lambda$. Thanks to the induction hypothesis, an application of \cref{Gluingaslowerbound}
to the (pairwise disjoint) clopen sets $(N_{s_n})_n$ allows then to conclude.

Consider the tree $T=\set{s\in\N^{<\N}}[\CB(g\corestr{N_s})=\lambda]$, notice that $T\neq\emptyset$ because it contains at least the empty sequence.
If $[T]$ is infinite then an application of \cref{InfiniteEmbedOmega} allows to find the desired sequence, so we can suppose that $[T]$ is finite.
Let $F$ be the set of $\sqsubset$-minimal elements of $\N^{<\N}\setminus T$. Then $\set{\CB(g\corestr{N_s})}[ s\in F]$ is a subset of $\lambda$ and we claim that it is cofinal in $\lambda$, which allows us to find the desired sequence. Towards a contradiction assume that for some $\beta<\lambda$ we have $\CB(g\corestr{N_s})<\beta$ for all $s\in F$. Then, by \cref{CBbasics0}~\cref{CBbasicsfromJSL2},  $\CB_{\beta}(g)\cap g^{-1}(N_s)=\emptyset$ for all $s\in F$ and so $\CB_{\beta}(g)\subseteq g^{-1}([T])$. But as $[T]$ is finite, we have $\CB_{\beta+1}(g)=\empty$ and so $\CB(g)\leq \beta+1$, a contradiction.
\end{proof}

\begin{remark}
Note that we actually have $\Maximalfct{n}\leq \Minimalfct{2n}$ for all $n\in\N$ and so the General Structure~\cref{JSLgeneralstructure} for functions of finite $\CB$-ranks is slightly stronger: for every $f,g$ with finite $\CB$-rank, $2\CB(f)\leq \CB(g)$ implies $f\leq g$.
\end{remark}

As a consequence we can finally prove that \cref{Introthm:LevelsAreFinitelyGenerated} implies \cref{Introthm:BQO,Introthm:BQOonScat}.
When a \bqo{} is also a partial order we call it a \emph{better partial order}.
We will use that a sum of \bqo{} indexed by a better partial order is itself still a \bqo{} (\cite[(9.14)]{simpsonbqo}).
The following result has the exact same proof as \cite[Theorem 5.4]{carroy2013quasi}.
\index[IdxNotion]{better-quasi-order (\bqo{})}

\begin{proposition}\label{FGgivesBQO_2}
If $\sC_{\beta}$ is \bqo{} for all $\beta<\alpha$, then $\sC_{<\alpha}$ is \bqo{}.
In particular, \cref{Introthm:LevelsAreFinitelyGenerated} implies \cref{Introthm:BQO,Introthm:BQOonScat}.
\end{proposition}

\begin{proof}
Consider the partial order $(\N, \leq^\bullet)$, where 
\[
m\leq^\bullet n \Lglra m=n \text{ or }2m<n.
\]
It is immediate to see that $(\N,\leq^\bullet)$ is a partial order and a \wqo{} and it is  \bqo{} by \cite[Example 1.6]{CarroyYPFromWell}. 
Similarly, we define a partial order $\leq^\bullet$ on $\omega_1$ by $\alpha_0\leq^\bullet \alpha_1$ if and only if $\alpha_0=\alpha_1$ or $2\alpha_0<\alpha_1$. As any ordinal $\alpha$ can be uniquely written as $\alpha=\lambda+n$ with $\lambda$ limit or null and $n\in\N$ and $2\alpha=\lambda+2n$, we have: 
\[
\lambda_0+n_0 \leq^\bullet \lambda_1+n_1 \Lglra \lambda_0<\lambda_1\mbox{ or } \bigl(\lambda_0=\lambda_1\mbox{ and } n_0\leq^\bullet n_1 \bigr).
\]
As this partial order is the sum of the \bqo{} $(\N,\leq^\bullet)$ along the well-ordered set of countable limit ordinals, it is \bqo{}.
Finally, consider the sum $S_{<\alpha}=\sum_{\beta<\alpha} \sC_{\beta}$ of the levels $\sC_\beta$, quasi-ordered by continuous reducibility, along the better partial order $\leq^\bullet$. By our hypothesis, $\sC_\beta$ is \bqo{} for all $\beta<\alpha$, so $S_{<\alpha}$ is \bqo{} too. 
 
Now observe that by the General Structure \cref{JSLgeneralstructure}, the mapping $h:\sC_{<\alpha}\to S_{<\alpha}$ given by $h(f)= (\CB(f),f)$ is a co-homomorphism. Namely, we have $f\leq g$ whenever $h(f)\leq h(g)$ in $S$, which by definition is equivalent to 
\[
\CB(f)<^\bullet \CB(g) \text{ or } (\CB(f)=\CB(g) \text{ and } f\leq g).
\]
Therefore, it follows that continuous reducibility is \bqo{} on $\sC_{<\alpha}$.

In particular, assuming \cref{Introthm:LevelsAreFinitelyGenerated} holds, then $\sC_\alpha$ is \bqo{} for all $\alpha<\omega_1$ by \cref{SecondstepforBQOthm}.  So $\sC$ is \bqo{} (\cref{Introthm:BQOonScat}) and hence \cref{Introthm:BQO} follows by \cref{FirststepforBQOthm}. 

\end{proof}

The General Structure \cref{JSLgeneralstructure} has many interesting consequences, we point out two of them that will reveal useful in the sequel.

\begin{corollary}\label{ConsequencesGeneralStructureThm}
Let $\lambda$ be a limit ordinal.
\begin{enumerate}
\item If $(f_n)_{n\in\N}$ is in $\sC_{<\lambda}$, then $\pgl_nf_n\leq\Minimalfct{\lambda+1}$,
and if moreover $(\CB(f_n))_n$ is regular and $\sup_n(\CB(f_n))=\lambda$, then $\pgl_n(f_n)\equiv \Minimalfct{\lambda+1}$.
\item If $\CB(f)\geq \lambda+2$, then $\pgl\Maximalfct{\lambda}\leq f$,\label{ConsequencesGeneralStructureThm2}
\end{enumerate}
\end{corollary}

\begin{proof} 
Fix any increasing cofinal sequence $(\alpha_n)_{n\in\N}$ in $\lambda$. There is an increasing sequence $(k_n)_{n\in\N}\in\cN$ such that $2\CB(f_n)\leq\alpha_{k_n}$ for all $n$.
So by \cref{JSLgeneralstructure} we have $f_n\leq \Minimalfct{\alpha_{k_n}+1}$ and in turn $\pgl_nf_n\leq\Minimalfct{\lambda+1}$ by \cref{Pgluingasupperbound}.
If moreover $(\CB(f_n))_n$ is regular and $\sup_n(\CB(f_n))=\lambda$, then by \cref{CBrankofPgluingofregularsequence1} $\CB(\pgl_nf_n)=\lambda+1$,
so by \cref{Minfunctions} $\Minimalfct{\lambda+1}$, being a minimum, it reduces to $\pgl_nf_n$.

For the second item, by \cref{JSLgeneralstructure} we have $\Maximalfct{\lambda}\leq \Minimalfct{\lambda+1}$ and so $\pgl \Maximalfct{\lambda}\leq \pgl \Minimalfct{\lambda+1} =\Minimalfct{\lambda+2}$. Since $\CB(f)\geq\lambda+2$, we have $\Minimalfct{\lambda+2}\leq f$ by \cref{Minfunctions} and so $\pgl \Maximalfct{\lambda}\leq f$, as desired.

For the last point, suppose towards a contradiction that they are equivalent and let $\Minimalfct{\lambda+1}=\pgl_n\Minimalfct{\alpha_n+1}$ for $(\alpha_n)_n$ cofinal in $\lambda$. By \cref{Rigidityofthecocenter}, it follows that $\Maximalfct{\lambda}\leq \gl_{n<M} \Minimalfct{\alpha_n+1}$ for some $M\in \N$, but then $\CB(\Maximalfct{\lambda})=\lambda\leq \sup_{n<M}(\alpha_n +1)<\lambda$. 
\end{proof}



{\cref{LocallyConstantFunctions} gives us the structure of $\sC_1$ (in other words, locally constant functions), and
the General Structure \cref{JSLgeneralstructure} tells us that functions with same limit rank are all continuously equivalent.
The next cases are those of $\sC_2$, and $\sC_{\lambda+1}$ for limit $\lambda$.
Interestingly, these two cases have more in common than just being the next ones to deal with,
see \cref{cor:CenteredSucessor,FGatsuccessoroflimit,OptimalityOfGenerators,qu:isomorphismofthelevels}.

